%% ============================================================
%%  CHAPTER 10 — CONCLUSIONS AND FUTURE WORK
%% ============================================================
\chapter{Conclusions and Future Work}
\label{ch:conclusion}

\section{Summary of Contributions}

This thesis has presented the \VIBE{} framework---a scalable, fully differentiable software architecture designed to bridge the gap between high-fidelity electrochemical simulations and system-level battery pack analysis. Moving beyond traditional numerical implementations, this work delivered the following core contributions directly aligned with the research objectives:

\begin{enumerate}
  \item \textbf{Fully Differentiable Physics-Informed Simulator.} Integrated a reverse-mode automatic differentiation pipeline using PyTorch that completely bypasses traditional finite-difference approximations. This provides exact, machine-precision gradient computation for complex sensitivity analyses (e.g., SOC-dependent temperature-rise sensitivities) directly through the highly coupled Dual-Foil Node (DFN) PDE constraints.

  \item \textbf{Modular, Plug-and-Play PDE Architecture.} Engineered an object-oriented framework that distinctly decouples the mathematical formulation of battery physics from the underlying spatial discretisation. This architecture permits runtime swapping of numerical backends and the dynamic, on-demand activation of various degradation states (such as SEI growth mechanisms) without rewriting core solver logic, ensuring unparalleled research flexibility.

  \item \textbf{Scalable GPU-Accelerated Pack Execution.} Addressed the critical computational bottleneck of large-scale pack simulations by developing a highly parallelised solver. By leveraging GPU batching (\texttt{torch.vmap}) and an implicit Newton-Raphson scheme with a block-diagonal Jacobian formulation, the framework successfully scales full-physics DFN models to multi-cell topologies (e.g., 10s10p packs) within drastically reduced execution times.

  \item \textbf{Coupled Heterogeneity and Degradation Modelling.} Established a robust software linkage between electrical network topologies, 1D thermal propagation, and internal electrochemical states. This enabled the successful simulation of cell-to-cell parameter variations, accurately capturing how thermal imbalances and contact resistance variations propagate through the pack to drive accelerated, localised SEI degradation.

  \item \textbf{Highly Optimised Simulation Data Pipeline.} Re-architected the simulation output infrastructure to overcome the massive data overhead generated by prolonged cycling of large battery packs. By introducing a highly configurable \texttt{OutputSpec} system, centralised \texttt{ResultsProcessor}, and compressed binary storage, the framework achieved immense I/O efficiency, entirely eliminating the latency of conventional text-based logging.
\end{enumerate}

\section{Answers to Research Questions}

\begin{description}
  \item[RQ1: Can DFN-fidelity pack simulation be made computationally feasible beyond small cell counts?]
  Yes. The block-diagonal Jacobian structure, GPU batching, and implicit timestepping allow the framework to simulate 25-cell (5S5P) packs with full DFN physics within minutes on a modern GPU.

  \item[RQ2: Does pack heterogeneity create self-reinforcing degradation feedback?]
  Yes. The 5S5P heterogeneity study showed that a single branch with $10\times$ lower cooling and $20$--$30\%$ higher contact resistance accumulates $\approx 2\times$ more SEI thickness after 100 cycles compared to nominal branches, driven by the thermal-electrochemical coupling.

  \item[RQ3: Do different SEI mechanisms produce quantifiably different degradation trajectories?]
  Yes. The seven SEI models span a $>10\times$ range in capacity fade rate under identical conditions. The tunnelling model is most aggressive at early cycles; the solvent-diffusion model produces the characteristic parabolic $L_\text{SEI}(t)$ growth; reaction-limited growth is most temperature-sensitive.

  \item[RQ4: Does physics-based evaluation change the assessment of BMS strategies?]
  Yes. Because the thermal heterogeneity persists even with perfect voltage balancing, BMS strategies evaluated against ECM simulations may over-estimate their effectiveness. Active liquid cooling provides $>2\times$ better temperature uniformity than passive balancing alone.
\end{description}

\section{Limitations}

\begin{itemize}
  \item \textbf{Spatial Resolution in Thermal Modelling:} The current thermal model assumes a lumped parameter approach (single uniform temperature per cell). Capturing intra-cell thermal gradients, particularly for large-format prismatic or pouch cells, requires 2D or 3D thermal coupling, which is currently unsupported natively by the PDE compiler.
  \item \textbf{Circuit Topology Assumptions:} The current distribution solver relies on an idealised Kirchhoff network topology. Contact resistance asymmetry or complex custom busbar geometries are not automatically parsed and require manual augmentation of the pack adjacency matrix.
  \item \textbf{Solver Memory Scaling:} While the block-diagonal Jacobian with GPU batching provides excellent speedup, the memory footprint scales linearly with the number of cells. For pack sizes exceeding several thousand cells, GPU VRAM limitations become a bottleneck for single-device execution.
  \item \textbf{Adaptive Timestepping:} The implicit Newton solver currently relies on simple heuristics for timestepping. A robust, fully adaptive step-size controller based on local truncation error estimation would improve computational efficiency during highly stiff transient periods.
\end{itemize}

\section{Future Work}
\label{sec:future_work}

\begin{enumerate}
  \item \textbf{Multi-GPU Distributed Simulation:} Extend the framework to support multi-GPU parallelisation via distributed tensor operations (e.g., PyTorch DDP). This would allow partitioning the pack adjacency matrix across multiple devices, enabling the simulation of extremely large utility-scale battery systems.

  \item \textbf{Advanced Preconditioning:} Implement domain-specific preconditioners (such as block Jacobi or algebraic multigrid) to accelerate the convergence of the linear solvers within the Newton-Raphson iterations, particularly for highly heterogeneous pack configurations where the Jacobian becomes ill-conditioned.

  \item \textbf{Dynamic Adaptive Mesh Refinement (AMR):} Introduce AMR capabilities into the spatial discretisation backends. Dynamically allocating computational nodes to regions with high spatial gradients would significantly reduce the degrees of freedom while maintaining numerical accuracy.

  \item \textbf{End-to-End Parameter Optimisation Pipeline:} Fully exploit the reverse-mode automatic differentiation capabilities by building a gradient-based parameter estimation module. This would allow automated fitting of framework parameters to experimental pack-level charge/discharge data using modern optimisers like Adam or L-BFGS.

  \item \textbf{Machine Learning Surrogate Embedding:} Utilise the framework to train Physics-Informed Neural Networks (PINNs) or neural operators, which can then be embedded back into the simulation pipeline to replace computationally expensive sub-models, creating a hybrid mechanistic-data-driven simulator.
\end{enumerate}
